\par\noindent\textbf{(1)} 珠子在惯性系的速度平方为 a²(θ̇²+Ω²sin²θ)，势能取 U=−mga cosθ。因此 L=ma²θ̇²/2+ma²Ω²sin²θ/2+mga cosθ。由 Euler–Lagrange 得 θ̈=(Ω²cosθ−g/a)sinθ；必须包含保持圆环匀速旋转的外驱动。
\par\noindent\textbf{(2)} 静止满足 sinθ(Ω²cosθ−g/a)=0，故 θ=0、π；当 Ω²≥g/a 另有对称的 θ=±arccos[g/(aΩ²)]。阈值 Ω\_c=√(g/a) 处分叉从最低点连续生出；上方根只能在非零转速足够大时存在。
\par\noindent\textbf{(3)} 令有效势 V\_eff=−mga cosθ−ma²Ω²sin²θ/2，线化频率平方为 V\_eff''/(ma²)= (g/a)cosθ−Ω²cos2θ。最低点为 g/a−Ω²，仅阈值以下稳定；顶部恒为 −g/a−Ω² 不稳定；倾斜分支代入 cosθ=g/(aΩ²)，得到 ω²=Ω²−g²/(a²Ω²)>0。
\par\noindent\textbf{(4)} 在旋转系时间不显含，约化守恒量 J=ma²θ̇²/2+V\_eff(θ)；它不是惯性系机械能，因为电机可通过约束做功。Ω→0 时回到 θ̈=−(g/a)sinθ 与普通摆频率；临界 Ω=√(g/a) 二阶刚度为零，需要四阶展开才能判断软模。
